\scanpage{48}
\mspage{4)}
% archive sheet number: 47
On a un foncteur
\diagram{
\begin{tikzcd}[ampersand replacement=\&]
(*)\qquad \Gamma:\quad S_{*}{\unsure{\mathscr T|\mathcal T}}D_{\mathrm{top}}^{\pm}\arrow[r] \&
\deleted{\unsure{\mathrm{FMD}|\mathrm{GMD}}}(\mathrm{Grabipond})
\end{tikzcd}
}
qui associe à un objet \inserted{$(X,\widetilde{\Sigma},R)$} de $S_{*}\mathscr{T}\deleted{\unsure{D}}D_{\mathrm{top}}^{\pm}$ le graphe \inserted{MD} $(S,\vec A,\sigma,\deleted{\unsure{\pi}}o,\unsure{\gamma|g},\unsure{\rho|r})$ suivant :
\[
S=\pi_{0}(X\deleted{\unsure{/}}\setminus\widetilde{\Sigma})\simeq
\pi_{0}(\mathrm{Déc}(X,\widetilde{\Sigma}))
\]
\[
\begin{aligned}
\vec A={}&\text{ens des couples }(C,\deleted{\unsure{o}}\omega),\ \text{où }C\ \deleted{\unsure{\text{est un}}}\inserted{\in\pi_{0}(\widetilde{\Sigma})},\ \text{et où}\\
&\omega\in\mathrm{Or}(C)\inserted{\simeq\mathrm{rive}_{X}(C)},\ \text{soumise à la condition que si}\\
&C\in\pi_{0}(\partial X),\ \text{alors ce soit l'orientation induite}\\
&\text{i.e. la rive qui se trouve dans }X
\end{aligned}
\]
\[
\begin{aligned}
\sigma={}&\text{échange des rives (quand on a le droit d'échanger,}\\
&\text{sinon, i.e. si }C\in\pi_{0}(\partial X),\ \text{on pose }\sigma(a)=a\text{)}
\end{aligned}
\]
$\deleted{\unsure{o}}o$ : \deleted{\unsure{Détermine}} associe à $(C,\text{rive de }C)$ l'unique ``\unsure{face}'' $\in\pi_{0}(X\setminus\widetilde{\Sigma})$ qui contient cette rive
\diagram{
\begin{tikzcd}[ampersand replacement=\&]
\gamma:\ \underset{\substack{\Vert\\\pi_{0}(X\setminus\widetilde\Sigma)}}{S}\arrow[r] \& \mathbb N
\end{tikzcd}
}
- application \uline{genre} des composantes $X_{\unsure{i|1}}$.
\diagram{
\begin{tikzcd}[ampersand replacement=\&]
\deleted{\unsure{\rho}}r:\ A=\pi_{0}(\widetilde{\Sigma})\arrow[r] \& \mathbb N\ :\ C\longmapsto\mathrm{Card}(R\cap C).
\end{tikzcd}
}

Le foncteur $(*)$ \deleted{est ess. \uline{surjectif} et} \inserted{induit un \uline{iso} sur les} $\pi_{0}$ et \unsure{une|un} \unsure{surjectif|surject.} sur les $\pi_{1}$ (c'est un $0$-\uline{isomorphisme} \deleted{\unsure{)|?}}, en langage homotopique).

\deleted{\unsure{Le fonct.}}\quad

\uline{Remarques} Dans $S_{*}\mathscr{T}D_{\mathrm{top}}^{\pm}$ il y a une opération évidente de \uline{somme} \inserted{\uline{directe}}, dans \deleted{FMD}\inserted{Grabipond} aussi, et $(*)$ y commute.
